%% =================================================================
%% Beber, Lamon, Saveriano, Fontanelli, Palopoli.
%% "Autonomous Robotic Tissue Palpation and Abnormalities
%%  Characterisation via Ergodic Exploration."
%% IEEE Robotics and Automation Letters, preprint accepted Feb. 2026.
%% DOI: 10.1109/LRA.2026.3673907
%%
%% Reconstruction of page 5 as study notes.
%% Contents: Figure 2 (three synthetic stiffness distributions),
%% Table I (performance metrics: 50 runs, noise-free vs noisy),
%% prose for Simulation Setup + start of section IV-B.
%% No NUMBERED equations; one inline unnumbered BO acquisition function.
%% Preamble adds booktabs/multirow/array for Table I.
%% =================================================================

\documentclass[11pt]{article}

\usepackage[margin=1in]{geometry}
\usepackage{amsmath, amssymb}
\usepackage{mathrsfs}
\usepackage{bm}
\usepackage{xcolor}
\usepackage{fancyhdr}
\usepackage{url}
\usepackage{booktabs}
\usepackage{multirow}
\usepackage{array}

\pagestyle{fancy}
\fancyhf{}
\lhead{\small Beber et al.: Ergodic Palpation for Tissue Mapping}
\rhead{\small IEEE RA-L Preprint, Accepted Feb.\ 2026}
\cfoot{\small 5 $\to$ \arabic{page}}
\renewcommand{\headrulewidth}{0.4pt}

\setcounter{page}{1}
\setcounter{equation}{6}   % page 4 ended with eq. (6); no numbered eqs on page 5
\setlength{\parskip}{4pt}

\begin{document}

% ---------- Preprint notice strip ----------
\begin{center}
{\footnotesize\itshape
This article has been accepted for publication in IEEE Robotics and
Automation Letters. This is the author's version which has not been
fully edited and content may change prior to final publication.
Citation information: DOI 10.1109/LRA.2026.3673907\par}
\end{center}

\vspace{0.4em}
\hrule
\vspace{0.6em}

\noindent
{\footnotesize\scshape Beber et al.: Ergodic Palpation for Tissue Mapping \hfill 5}

\vspace{1.2em}

% ---------- Continuation of IV-A Setup and Baselines from page 4 ----------
\noindent
They were generated as mixtures of Gaussian components with varying
locations, covariances, and amplitudes, modelling localised stiff
inclusions (over 100~kPa in pathological areas) in a softer
background (few kilopascal in healthy regions) [29], similarly to
related research [15]--[17]. Three scenarios in a square domain
(with side 45~mm) were considered: (i) a single stiff region,
(ii) two stiff regions, and (iii) three distinct stiff areas, whose
distributions are illustrated in Fig.~2.

% ---------- Figure 2 placeholder ----------
\vspace{0.4em}
\begin{center}
\fbox{\parbox{0.92\textwidth}{\centering\vspace{1em}
\textbf{[Figure 2 --- placeholder]}\\[0.8em]
Three side-by-side 2D colour maps showing the synthetic ground-truth
elasticity distributions used as simulation targets, labelled
\textit{1 Stiff Area}, \textit{2 Stiff Areas}, \textit{3 Stiff
Areas}. Each panel plots $x_{1}$ horizontally and $x_{2}$ vertically
(both in mm, on a 0--45 grid). A shared vertical colour bar on the
right encodes the elastic modulus in kPa, with the background at a
low value (a few kPa, healthy tissue) and localised peaks above
100~kPa marking the stiff inclusions. The inclusion patterns are
Gaussian blobs of varying covariance and amplitude, placed at
different locations in each of the three panels.
\vspace{1em}}}
\end{center}

\vspace{0.3em}

\begin{center}
\textbf{Fig. 2.} Synthetic elasticity distributions for simulated
scenarios, from 1 to 3 stiffer regions. The colour scale refers to
the elastic modulus (in kPa).
\end{center}

\vspace{0.8em}

\noindent
These regions were first densely sampled with a grid-based approach
to generate the ground truth, used to compute the estimation and
segmentation errors. The simulation environment was implemented in
Python, with GPR and BO components built using the \texttt{PyTorch}
framework, thus ensuring efficient GPU-based inference and
optimisation during online execution. All simulations and
experiments were executed on a mini PC equipped with an Intel
i7-13700HX CPU (24 cores), 64~GB RAM, and an NVIDIA RTX 4070 GPU
with 8~GB of memory. The proposed framework runs in real time, with
HEDAC trajectory updates and EKF-based viscoelastic estimation both
requiring less than 1~ms per control loop iteration at 100~Hz. GPR
updates are performed asynchronously at a lower frequency (1~Hz),
with a maximum observed update time of 105~ms. Elasticity samples
were collected at a frequency of 4~Hz and a maximum velocity of
1~cm/s. In this configuration, the maximum distance between two
consecutive sampled points equals half of the indenter radius, which
is set here to 2.5~mm. These simulations were conducted to evaluate
the performance of the proposed ergodic search strategy
(Sec.~IV-B) across the three distribution scenarios.

\vspace{0.4em}

\noindent
We compared ergodic search with two variants of Bayesian
Optimisation (BO), which is the standard baseline for robotic
palpation (Sec.~IV-C). BO models an expensive-to-evaluate objective
function $f(x)$ using a Gaussian Process (GP) and selects sampling
locations by maximising an acquisition function that balances
exploration and exploitation. Given data
$\mathcal{D} = (x_{i}, y_{i})_{i=1}^{n}$, the GP provides a
posterior predictive distribution
$f(x) \sim \mathcal{N}(\mu(x),\,\sigma^{2}(x))$. As an acquisition
function, the Expected Improvement (EI) [15], [16] is employed,
$a_{\mathrm{EI}}(x) = (f_{\min} - \mu(x))\,\Phi(z) + \sigma(x)\,
\phi(z)$, where $f_{\min}$ is the minimum observed value and
$z = \dfrac{f_{\min} - \mu(x)}{\sigma(x)}$. This variant is denoted
BO-EI. However, the standard point-to-point implementation of BO has
limitations for robotic palpation. Due to safety constraints,
especially in medical settings, the robot must follow smooth and
predictable trajectories and cannot move rapidly between distant
sampling points as typically generated by the BO. Moreover, each
probing action can take up to 5~s due to indentation and
viscoelastic estimation time. To ensure a fair comparison with the
proposed ergodic planner, we adapted BO to perform continuous
palpation, which characterises our ergodic search. This prevents BO
from penalising non-sampled regions during motion. Specifically,
additional samples are collected every 2.5~mm (half the indenter
radius) along the path between BO-selected targets, enabling BO to
exploit intermediate data, similarly to the ergodic approach
(variant termed BO-EIS).

% ---------- Table I ----------
\vspace{0.8em}

\noindent\textbf{Table I}\\[0.2em]
{\small Performance metrics across 50 runs with noise-free elasticity
measurements (w/o) and with white noise $\mathcal{N}(0,\,6.25\,
\mathrm{kPa}^{2})$ (w/). NR is the number of stiff regions, DR the
detection rate, $l$ the lenght of the trajectory, and RMSE is the
root mean square error.}

\vspace{0.3em}

\begin{center}
\small
\begin{tabular}{c c c c c}
\toprule
NR & DR (\%) $\uparrow$ & $l$ [mm] & Time [s] $\downarrow$ & RMSE [kPa] $\downarrow$ \\
\midrule
\multicolumn{5}{c}{\textit{w/o Measurement Noise}} \\
\midrule
1 & 100\,\% & 664 & 70 & 0.37 \\
2 & 100\,\% & 660 & 69 & 0.97 \\
3 & 99\,\%  & 665 & 70 & 1.96 \\
\midrule
\multicolumn{5}{c}{\textit{w/\phantom{o} Measurement Noise}} \\
\midrule
1 & 100\,\% & 633 & 66 & 1.57 \\
2 & 99\,\%  & 630 & 65 & 2.17 \\
3 & 97\,\%  & 580 & 60 & 3.14 \\
\bottomrule
\end{tabular}
\end{center}

\vspace{0.6em}

\noindent
To assess the spatial accuracy of the reconstructed stiffness maps,
we applied a segmentation pipeline to identify distinct stiff
regions. Based on the elasticity estimates in Sec.~IV-B and IV-C,
KMeans clustering (\texttt{SciPy}) was used to separate stiffer
regions from the surrounding soft tissue. Unlike fixed thresholding
approaches [12], this method does not require manual parameter
tuning and adapts to the distribution of estimated stiffness values.
The number of detected regions was obtained by counting connected
components in the stiff cluster. In Sec.~IV-D, the segmented regions
are converted into polygonal boundaries using \texttt{Shapely} and
compared against ground truth using standard segmentation metrics.

\vspace{0.6em}

% ---------- B. Robustness of Ergodic Exploration ----------
\noindent\textit{B.~Robustness of Ergodic Exploration}

\vspace{0.2em}

\noindent
In this section, the performance of the ergodic controller is tested
with different numbers of stiffer regions. To assess the robustness
of the algorithm, for each map, the simulation is executed 50 times
with a random starting position, while three performance metrics
were considered: 1. the detection rate (DR), i.e., the percentage of
trials in which the algorithm correctly estimates the number and
position of the stiff regions present in the ground truth (higher is
better); 2. the elapsed time (lower is better); 3. the root mean
square error (RMSE, lower is better). For completeness, we report
also the trajectory length $l$. Each exploration terminates when the
ergodic metric $\alpha$ reaches 0.4 (set experimentally). Two
experimental conditions were examined: one involved noise-free
elasticity measurements and another incorporated noisy measurements
$\mathcal{N}(0,\,6.25\,\mathrm{kPa}^{2})$, where the variance equals
the noise of the elasticity estimation in the worst case. The
results of the study are reported in Tab.~I. Overall, the algorithm
exhibited comparable performance in terms of DR and execution

% ---------- Bottom license strip ----------
\vfill
\hrule
\vspace{0.3em}
\begin{center}
{\footnotesize\itshape
This work is licensed under a Creative Commons Attribution 4.0
License. For more information, see
\url{https://creativecommons.org/licenses/by/4.0/}\par}
\end{center}

\end{document}
